\documentclass[10pt,a4paper]{amsart}
\usepackage[margin=19mm]{geometry}
\usepackage{amsmath,amssymb,mathtools}
\usepackage[scr=rsfs,cal=euler]{mathalfa}
\usepackage{xfrac}
\usepackage[unicode,hidelinks,bookmarks=false]{hyperref}
\newcommand{\ie}{i.e.\ }
\newcommand{\cf}{cf.\ }
\newcommand{\resp}{respectively\ }
\newcommand{\Subsection}{\subsection}
\providecommand{\nobreakdash}{\nobreak\hbox{-}}
\setlength{\emergencystretch}{2em}
\multlinegap0pt
\usepackage{amssymb}
\newtheorem{lem}{}
\newtheorem{thm}{}
\title{On the identifiability of Dirichlet mixture models}
\author{Hien Duy Nguyen, Mayetri Gupta}
\date{Source: arXiv:2603.21914v1 (CC BY 4.0)}
\begin{document}
\maketitle

\noindent\emph{Source and licence.} On the identifiability of Dirichlet mixture models. Version arXiv:2603.21914v1 (CC BY 4.0), distributed by arXiv under the Creative Commons Attribution 4.0 International licence, http://creativecommons.org/licenses/by/4.0/, as recorded in the arXiv licence metadata for this exact version and shown on the abstract page https://arxiv.org/abs/2603.21914v1. Full licence notice: \texttt{licenses/arxiv-2603.21914-CC-BY-4.0.txt}.\par\smallskip

\noindent\textit{Editorial scope.} Counted unit: the article's thm thm:linind\_FD (source0.tex lines 619--623). The article's complete original proof of it is delivered with it (source0.tex lines 625--788, one \texttt{proof} environment of 164 lines). No mathematics is written, repaired or re-derived: the statement and the proof are the article's own lines, in the article's own notation, and no retained line is substituted or re-flowed.  Retained uncounted as the source-local context the proof needs: lines 298--307; lines 545--548; lines 555--564; lines 587--590; lines 613--615.  Nothing was omitted from any retained span, so the package carries no omission record.  No retained line contains a citation, so the wrapper carries no bibliography and no bibliography entry.  No retained line contains a figure environment, an image-inclusion command or drawing code, so no image is delivered and the package declares no asset dependency.  Editorial formatting only, in addition: an AMS \texttt{amsart} preamble that restates the article's own declarations and macros used by the retained lines, namely the environments \texttt{lem}, \texttt{thm} on separate counters, together with the standard packages those lines need.  The retained environments print in this document as Lemma 1, Theorem 1 in order of appearance, whereas the article prints the same items with the numbering of its own full document.  The complete native source of the article as distributed in the arXiv e-print for this exact version is bundled unchanged as \texttt{sources/196983/source0.tex}.
\begin{lem}
\label{lem:linind_ident} Let $\Theta\subseteq\mathbb{R}_{>0}^{J}$.
Suppose the family $\{f_{\alpha}:\alpha\in\Theta\}$ is linearly independent
in the sense that for any distinct $\alpha^{(1)},\dots,\alpha^{(K)}\in\Theta$,
$K\in\mathbb{N}$, 
\[
\sum_{k=1}^{K}c_{k}f_{\alpha^{(k)}}(x)=0\ \text{a.e. on }\Delta_{J-1}^{\circ}\quad\Longrightarrow\quad c_{1}=\cdots=c_{K}=0.
\]
Then $\mathcal{F}\left(\Theta\right)$ is identifiable. 
\end{lem}

\begin{lem}
\label{lem:disjoint} Let $u,u'\in(-1,0)^{d}$ and $m,m'\in\mathbb{N}_{0}^{d}$.
If $u+m=u'+m'$, then $u=u'$ and $m=m'$. 
\end{lem}

\begin{lem}
\label{lem:dirichlet_series} Let $\{\lambda_{n}\}_{n\ge1}$ be a
strictly increasing sequence of real numbers with $\lambda_{n}\to\infty$.
Suppose $\sum_{n\ge1}a_{n}\mathrm{e}^{-\lambda_{n}s}$ converges absolutely
for all $s\ge s_{0}$ and 
\[
\sum_{n\ge1}a_{n}\mathrm{e}^{-\lambda_{n}s}=0\qquad\forall s\ge s_{0}.
\]
Then $a_{n}=0$ for all $n$. 
\end{lem}

\begin{lem}
\label{lem:chart_jac} The Jacobian determinant of $y\mapsto(x_{1}(y),\dots,x_{d}(y))$
is $(1+Y)^{-J}$. 
\end{lem}

\begin{equation}
h_{\alpha}(y)=\frac{\Gamma(\alpha_{+})}{\prod_{j=1}^{J}\Gamma(\alpha_{j})}\frac{\prod_{j=1}^{d}y_{j}^{\alpha_{j}-1}}{(1+Y)^{\alpha_{+}}}.\label{eq:transported_kernel}
\end{equation}

\begin{thm}
\label{thm:linind_FD} The family $\{f_{\alpha}:\alpha\in\Theta_{\mathrm{FD}}\}$
is linearly independent. Consequently, the finite Dirichlet mixture
class $\mathcal{F}\left(\Theta_{\mathrm{FD}}\right)$ is identifiable. 
\end{thm}

\begin{proof}
By Lemma~\ref{lem:linind_ident}, it is enough to prove linear independence.
Let $\alpha^{(1)},\dots,\alpha^{(K)}\in\Theta_{\mathrm{FD}}$ be distinct
and suppose that 
\[
\sum_{k=1}^{K}c_{k}f_{\alpha^{(k)}}(x)=0\qquad\text{for a.e. }x\in\Delta_{J-1}^{\circ}.
\]
The left-hand side is continuous on $\Delta_{J-1}^{\circ}$, so the
identity in fact holds for all $x\in\Delta_{J-1}^{\circ}$.

Transport this identity by the bijection $y\mapsto\varphi(y)$. Multiplying
by the Jacobian determinant from Lemma~\ref{lem:chart_jac}, which
is strictly positive and does not depend on $\alpha$, yields 
\begin{equation}
\sum_{k=1}^{K}c_{k}h_{\alpha^{(k)}}(y)=0\qquad\text{for all }y\in\mathbb{R}_{>0}^{d}.\label{eq:FD_transport_identity}
\end{equation}
Fix the open neighbourhood 
\[
U=\{y\in\mathbb{R}_{>0}^{d}:Y<1/2\}.
\]
For $v>0$ and $y\in U$, the binomial series gives 
\begin{equation}
(1+Y)^{-v}=\sum_{n=0}^{\infty}(-1)^{n}\frac{(v)_{n}}{n!}Y^{n},\qquad(v)_{n}=v(v+1)\cdots(v+n-1),\label{eq:FD_binomial}
\end{equation}
and for each $n\in\mathbb{N}_{0}$ the multinomial theorem gives 
\begin{equation}
Y^{n}=(y_{1}+\cdots+y_{d})^{n}=\sum_{\substack{m\in\mathbb{N}_{0}^{d}\\
|m|=n
}
}\frac{n!}{m!}y^{m}.\label{eq:FD_multinomial}
\end{equation}
Now fix $\alpha\in\Theta_{\mathrm{FD}}$ and set 
\[
u(\alpha)=(\alpha_{1}-1,\dots,\alpha_{d}-1)\in(-1,0)^{d},\qquad v(\alpha)=\alpha_{+}.
\]
Substituting \eqref{eq:FD_binomial} and \eqref{eq:FD_multinomial}
into \eqref{eq:transported_kernel} yields, for $y\in U$, 
\begin{equation}
h_{\alpha}(y)=\frac{\Gamma(v(\alpha))}{\prod_{j=1}^{J}\Gamma(\alpha_{j})}\sum_{m\in\mathbb{N}_{0}^{d}}(-1)^{|m|}\frac{(v(\alpha))_{|m|}}{m!}y^{u(\alpha)+m}.\label{eq:FD_series_expansion}
\end{equation}
For each fixed $y\in U$, the series in \eqref{eq:FD_series_expansion}
is absolutely convergent. Indeed, using \eqref{eq:FD_multinomial},
\[
\sum_{m\in\mathbb{N}_{0}^{d}}\frac{(v(\alpha))_{|m|}}{m!}y^{m}=\sum_{n=0}^{\infty}(v(\alpha))_{n}\sum_{\substack{m\in\mathbb{N}_{0}^{d}\\
|m|=n
}
}\frac{y^{m}}{m!}=\sum_{n=0}^{\infty}(v(\alpha))_{n}\frac{Y^{n}}{n!},
\]
and the last series converges absolutely because $Y<1$. Since $y^{u(\alpha)}=\prod_{j=1}^{d}y_{j}^{\alpha_{j}-1}$
is finite for every $y\in U$, \eqref{eq:FD_series_expansion} is
pointwise absolutely convergent on $U$.

Substituting \eqref{eq:FD_series_expansion} into \eqref{eq:FD_transport_identity}
is legitimate for each $y\in U$, since we are summing only finitely
many pointwise absolutely convergent series. Moreover, if
\[
u(\alpha^{(k)})+m=u(\alpha^{(k')})+m',
\]
then
\[
u(\alpha^{(k)})-u(\alpha^{(k')})=m'-m\in\mathbb{Z}^{d},
\]
so Lemma~\ref{lem:disjoint} implies that $u(\alpha^{(k)})=u(\alpha^{(k')})$,
and then necessarily $m=m'$. Thus coincidences among exponent vectors
can occur only among terms having the same value of $u(\alpha)$, and
regrouping first by the distinct values of $u(\alpha)$ is unambiguous.
Let $u_{1},\dots,u_{L}\in(-1,0)^{d}$ be the distinct values taken
by $u(\alpha^{(k)})$, and let $\mathcal{G}_{g}$ be the set of indices
$k$ such that $u(\alpha^{(k)})=u_{g}$. For $k\in\mathcal{G}_{g}$,
write 
\[
v_{k}=\alpha_{+}^{(k)},\qquad d_{k}=c_{k}\frac{\Gamma(v_{k})}{\prod_{j=1}^{J}\Gamma(\alpha_{j}^{(k)})}.
\]
Then we obtain 
\begin{equation}
\sum_{g=1}^{L}\sum_{m\in\mathbb{N}_{0}^{d}}A_{g,m}y^{u_{g}+m}=0\qquad\text{for all }y\in U,\label{eq:FD_regrouped}
\end{equation}
where 
\begin{equation}
A_{g,m}=\sum_{k\in\mathcal{G}_{g}}d_{k}(-1)^{|m|}\frac{(v_{k})_{|m|}}{m!}.\label{eq:FD_Agm}
\end{equation}
By Lemma~\ref{lem:disjoint}, all exponent vectors $u_{g}+m$ appearing
in \eqref{eq:FD_regrouped} are distinct.

Let 
\[
W=\{u_{g}+m:1\le g\le L,\ m\in\mathbb{N}_{0}^{d}\}.
\]
This is a countable subset of $\mathbb{R}^{d}$. For any two distinct
vectors $w,w'\in W$, the constraint $\langle w,\lambda\rangle=\langle w',\lambda\rangle$
defines a hyperplane in $\mathbb{R}^{d}$. The union of these hyperplanes
over all pairs $w\neq w'$ is therefore a countable union of Lebesgue-null
sets, hence cannot cover the positive orthant $\mathbb{R}_{>0}^{d}$.
Choose $\lambda\in\mathbb{R}_{>0}^{d}$ outside this union. Then the
scalar products $\langle w,\lambda\rangle$ are pairwise distinct
for $w\in W$.

For $s\ge0$, define 
\[
y(s)=(\mathrm{e}^{-\lambda_{1}s},\dots,\mathrm{e}^{-\lambda_{d}s}).
\]
Since each $\lambda_{j}>0$, we have $Y(s)=\sum_{j=1}^{d}\mathrm{e}^{-\lambda_{j}s}\to0$,
so $y(s)\in U$ for all $s\ge s_{0}$ for some $s_{0}$. Substituting
$y=y(s)$ into \eqref{eq:FD_regrouped} gives 
\begin{equation}
\sum_{g=1}^{L}\sum_{m\in\mathbb{N}_{0}^{d}}A_{g,m}\mathrm{e}^{-\mu_{g,m}s}=0\qquad\text{for all }s\ge s_{0},\label{eq:FD_dirichlet_series}
\end{equation}
where $\mu_{g,m}=\langle u_{g}+m,\lambda\rangle$. For each fixed
$s\ge s_{0}$, the series in \eqref{eq:FD_dirichlet_series} is absolutely
convergent because \eqref{eq:FD_regrouped} is absolutely convergent
at $y(s)\in U$. Moreover, the exponents $\mu_{g,m}$ are all distinct,
and they can be arranged into a strictly increasing sequence diverging
to $\infty$. Indeed, if $\lambda_{\min}=\min_{1\le j\le d}\lambda_{j}$
and $C=\min_{1\le g\le L}\langle u_{g},\lambda\rangle$, then 
\[
\mu_{g,m}\ge C+\langle m,\lambda\rangle\ge C+\lambda_{\min}|m|,
\]
so only finitely many pairs $(g,m)$ satisfy $\mu_{g,m}\le M$ for
a given $M$. Lemma~\ref{lem:dirichlet_series} therefore applies
to \eqref{eq:FD_dirichlet_series} and yields 
\begin{equation}
A_{g,m}=0\qquad\text{for all }g\in\{1,\dots,L\}\text{ and all }m\in\mathbb{N}_{0}^{d}.\label{eq:FD_Agm_zero}
\end{equation}

Now fix $g\in\{1,\dots,L\}$. Combining \eqref{eq:FD_Agm} and \eqref{eq:FD_Agm_zero}
gives 
\begin{equation}
\sum_{k\in\mathcal{G}_{g}}d_{k}(v_{k})_{|m|}=0\qquad\text{for all }m\in\mathbb{N}_{0}^{d}.\label{eq:FD_group_pochhammer}
\end{equation}
In particular, taking $m=(n,0,\dots,0)$ gives 
\begin{equation}
\sum_{k\in\mathcal{G}_{g}}d_{k}(v_{k})_{n}=0\qquad\text{for all }n\in\mathbb{N}_{0}.\label{eq:FD_group_pochhammer_n}
\end{equation}
Multiply \eqref{eq:FD_group_pochhammer_n} by $z^{n}/n!$ and sum
over $n\ge0$. Since $\sum_{n\ge0}(v)_{n}z^{n}/n!=(1-z)^{-v}$ for
$|z|<1$, we obtain 
\begin{equation}
\sum_{k\in\mathcal{G}_{g}}d_{k}(1-z)^{-v_{k}}=0\qquad\text{for all }z\in(0,1).\label{eq:FD_group_binomial}
\end{equation}
Now set $r=-\log(1-z)\in\mathbb{R}_{>0}$. Then \eqref{eq:FD_group_binomial}
becomes 
\begin{equation}
\sum_{k\in\mathcal{G}_{g}}d_{k}\mathrm{e}^{v_{k}r}=0\qquad\text{for all }r>0.\label{eq:FD_group_exponential}
\end{equation}
The scalars $\{v_{k}:k\in\mathcal{G}_{g}\}$ are distinct. Indeed,
within a fixed group $\mathcal{G}_{g}$ the first $d=J-1$ coordinates
of $\alpha^{(k)}$ are fixed, so equality $v_{k}=v_{k'}$ would also
force 
\[
\alpha_{J}^{(k)}=v_{k}-\sum_{j=1}^{d}\alpha_{j}^{(k)}=v_{k'}-\sum_{j=1}^{d}\alpha_{j}^{(k')}=\alpha_{J}^{(k')},
\]
and hence $\alpha^{(k)}=\alpha^{(k')}$, contrary to the distinctness
of the parameters. Writing $\mathcal{G}_{g}=\{k_{1},\dots,k_{r}\}$
and ordering so that $v_{k_{1}}<\cdots<v_{k_{r}}$, we divide \eqref{eq:FD_group_exponential}
by $\mathrm{e}^{v_{k_{r}}r}$ and let $r\to\infty$. This shows that
the coefficient of $\mathrm{e}^{v_{k_{r}}r}$ is $0$. Repeating the
argument gives $d_{k}=0$ for every $k\in\mathcal{G}_{g}$.

Finally, the factor $\Gamma(v_{k})/\prod_{j=1}^{J}\Gamma(\alpha_{j}^{(k)})$
is nonzero, so $d_{k}=0$ implies $c_{k}=0$ for all $k\in\mathcal{G}_{g}$.
Since $g$ was arbitrary, we conclude that $c_{1}=\cdots=c_{K}=0$.
This proves linear independence, and the identifiability statement
then follows from Lemma~\ref{lem:linind_ident}. 
\end{proof}

\end{document}
